\section{Method}
\label{sec:method}

\subsection{Event-guided assessment framework}
\label{ssec:framework}

Our goal is to assess EASE needle-handling sub-skills from a suturing video
alone. Each sub-skill, however, refers to a different part of a stitch. Needle
Repositions concerns how the needle is prepared for the next stitch, whereas
Needle Hold Ratio, Needle Hold Angle, and Depth of Needle Hold concern how the
needle is held when it is driven into the tissue. Scoring a whole video, or a
whole stitch, therefore mixes in frames that are irrelevant to the sub-skill.
We instead detect events that delimit each stitch, and assess each sub-skill
only within the interval it refers to.

A stitch generally proceeds as follows. After the previous needle has left the
tissue, the surgeon passes the needle between the two instruments to prepare
for the next stitch. The last of these repositions fixes the hold, which is
kept until the needle enters the tissue. We detect the five types of events
in Table~\ref{tab:events}, each an instant in the video. A stitch ends at a
withdrawal end, and its entry $e_i$ is the last entry before it; a retraction
does not end the stitch.

\begin{table}[htb]
\centering
\small
\caption{Event types.}
\label{tab:events}
\begin{tabular}{@{}ll@{}}
\hline
Event & Definition \\
\hline
Reposition start & Both instruments grasp the needle outside tissue \\
Reposition end & One instrument releases the needle \\
Entry & The needle tip enters the tissue \\
Withdrawal end & The needle leaves the tissue on the far side \\
Retraction end & The needle is pulled back out on the entry side \\
\hline
\end{tabular}
\end{table}

Two observation windows follow from these events (Fig.~\ref{fig:overview}). The reposition
window $R_i=(w_{i-1},e_i]$ spans from the previous withdrawal end $w_{i-1}$
(the start of the video for the first stitch) to the entry. The hold window
$W_i=[a_i,e_i]$ starts at the last reposition end $a_i$ in $R_i$, or at
$w_{i-1}$ if there is none. Since the hold does not change until entry, the
hold sub-skills can be measured from any frame in $W_i$. Following EASE,
Needle Repositions is scored from the number $n_i$ of reposition starts in
$R_i$: 3, 2, and 1 for $n_i\leq2$, $n_i=3$, and $n_i\geq4$, respectively.

\subsection{Per-frame observations}
\label{ssec:observations}

To measure the hold geometrically, each frame in the hold window is given
keypoints and a needle mask: the needle tip $P$ and tail $T$, the grasp point
$G$ where the jaws hold the needle, the two jaw tips with midpoint $J$, the
jaw base $H$, and the needle mask $M$. The framework does not prescribe how
these are obtained; a detector and a segmentation model trained on local
annotations would serve equally well. In this work, both the events and the
per-frame observations are obtained from a general-purpose VLM given a few
examples, without any training (Sec.~\ref{ssec:implementation}). The
framework can thus be set up without task-specific training, which should ease
its transfer to other suturing videos.

\subsection{Needle Hold Ratio}
\label{ssec:ratio}

\placeholder{Shared circular arc and the grasp position $s_i$.}

\subsection{Needle Hold Angle}
\label{ssec:angle}

\placeholder{Angle lifted to 3D with depth priors.}

\subsection{Depth of Needle Hold}
\label{ssec:depth}

\placeholder{Depth from the longest-jaw frame.}
